% Abstract written strictly from SUPPORTED claims verified in results/v3/.
% Every number is commented with its exact source file.
% Word count: ~165 words (target: <= 200 words).
\begin{abstract}
Text-proxy multimodal retrieval systems substitute expensive raw pixel models with text transcripts generated by optical character recognition (OCR) or image captioning. We evaluate when text proxies suffice across 12 operational regimes spanning sparse, dense, and visual retrievers on 626 receipt documents and 1,000 MS-COCO images.
% Source: results/v3/headline_hypothesis_tests.csv
We discover that semantic query mismatch induces a far more severe performance drop ($\Delta = 0.3782$ MRR) than clean OCR corruption ($\Delta = 0.0201$, $p = 10^{-5}$) or optical blur ($\Delta = 0.1600$, $p = 10^{-5}$) with non-overlapping 95\% confidence intervals.
% Source: results/v3/headline_regime_grid_v3.csv
Under verbatim queries with clean OCR, tuned BM25 achieves 0.8371 MRR, whereas dense bi-encoders struggle (MiniLM: 0.2056) due to WordPiece token fragmentation (+122\%) rather than sequence truncation (only 0.64\% receipts exceed 512 tokens).
% Source: results/v3/char_vs_word_bm25.csv
Character $n$-grams defend strictly against synthetic noise (+0.1008 MRR at 20\% noise), whereas word BM25 is superior under optical OCR (+0.0996 clean, +0.0777 blur; $p < 0.001$).
% Source: results/v3/coco_native_visual_v3.csv, results/v3/coco_hypothesis_test_v3.csv
On MS-COCO, native visual models (SigLIP: 0.6071 MRR; CLIP: 0.4709) decisively outperform caption proxies (MiniLM: 0.3728; $p_{\text{Holm}} = 8 \times 10^{-5}$).
% Source: results/v3/multihop_paired_results.csv
Finally, multi-hop generative bridging adds no detectable retrieval lift over single-stage hybrid retrieval (+1 net hit / 50; exact McNemar $p = 1.000$).
\end{abstract}
